Image GPTは、GPT系のautoregressive modelingを画像へ持ち込み、pixelをsequenceとして扱う研究方向を示した。languageとimageのdata structureは異なるが、前のtokenから次を予測するという一般的なsequence modelingの考え方をvisual domainへ適用した点に意味がある。2020年には、Transformer的な生成原理がtext以外でもどこまで成立するかが明示的な研究テーマになっていた。\autocite{src-d46ff01d5c8d}


Taming Transformersは、高解像度画像をpixel-levelの長いsequenceとして直接modelingすると計算量が大きくなる問題に対し、CNNを使ってcontext-richな離散visual vocabularyを学習し、そのcompositionをTransformerでmodelingする構成を示した。ここでは『画像を何にtokenizeするか』がarchitectureの中心になる。後年のlatent / discrete representation系へつながる設計空間を、2020年末の一次資料として確認できる。\autocite{src-03b656c5d291}


この系譜はDDPMなどのdiffusionと混同しない方がよい。Image GPTやTaming Transformersでは、離散的またはsequence-likeなrepresentationを順次modelingすることが中心であるのに対し、diffusion系ではnoiseを加えるforward processとdenoising / score estimationを通じてsampleを構成する。2020年のgenerative mediaは、一つの勝者へ収束する前に複数の生成原理が並立していた。\autocite{src-03b656c5d291,src-d46ff01d5c8d}


Image GPTからTaming Transformersまでを見ると、languageで強かったautoregressive sequence modelingをそのまま他modalitiesへ延長するだけでなく、画像側のrepresentationを変換してTransformerが扱いやすいsymbol列へ落とす方向も現れた。これは『生成AI=language model』という後年の単純化では捉えにくい2020年の分岐であり、翌年以降のmultimodal / generative-media研究を理解する前提になる。\autocite{src-03b656c5d291,src-d46ff01d5c8d}


